%% Roteiro — Entrega 1 — 5 minutos
%% pdflatex roteiro_entrega1.tex
\documentclass[12pt,a4paper]{article}

\usepackage[utf8]{inputenc}
\usepackage[T1]{fontenc}
\usepackage[brazil]{babel}
\usepackage[top=1.6cm,bottom=1.6cm,left=1.8cm,right=1.8cm,headheight=15pt]{geometry}
\usepackage[most]{tcolorbox}
\usepackage{tabularx}
\usepackage{booktabs}
\usepackage{enumitem}
\usepackage{fancyhdr}
\usepackage{parskip}
\usepackage{xcolor}
\usepackage{titlesec}

\definecolor{cfala}{RGB}{25,90,170}
\definecolor{cacao}{RGB}{15,130,55}
\definecolor{catencao}{RGB}{190,70,15}
\definecolor{ctempo}{RGB}{110,50,155}
\definecolor{bgfala}{RGB}{232,242,255}
\definecolor{bgacao}{RGB}{232,250,238}
\definecolor{bgatencao}{RGB}{255,240,225}
\definecolor{cinzabg}{RGB}{248,248,248}

\tcbset{base/.style={arc=3pt,boxrule=1.5pt,left=8pt,right=8pt,top=5pt,bottom=5pt,
  before upper={\parindent0pt}}}

\newtcolorbox{falabox}{base,colback=bgfala,colframe=cfala,
  title={\textbf{\textcolor{white}{FALA}}},coltitle=white,
  attach boxed title to top left={yshift=-2mm,xshift=6pt},
  boxed title style={colback=cfala,arc=2pt,boxrule=0pt}}

\newtcolorbox{acaobox}{base,colback=bgacao,colframe=cacao,
  title={\textbf{\textcolor{white}{AÇÃO}}},coltitle=white,
  attach boxed title to top left={yshift=-2mm,xshift=6pt},
  boxed title style={colback=cacao,arc=2pt,boxrule=0pt}}

\newtcolorbox{alertabox}{base,colback=bgatencao,colframe=catencao,
  title={\textbf{\textcolor{white}{ATENÇÃO}}},coltitle=white,
  attach boxed title to top left={yshift=-2mm,xshift=6pt},
  boxed title style={colback=catencao,arc=2pt,boxrule=0pt}}

\newcommand{\cmd}[1]{\texttt{\textbf{#1}}}

\newcommand{\acaoinline}[1]{%
  \par\vspace{3pt}%
  \noindent\colorbox{cacao!18}{%
    \parbox{\dimexpr\linewidth-2\fboxsep}{%
      \small\textcolor{cacao}{\textbf{[FAZER:}} #1\textcolor{cacao}{\textbf{]}}}}%
  \par\vspace{3pt}}

\newcommand{\tempo}[2]{%
  \vspace{5pt}
  \begin{tcolorbox}[base,colback=ctempo!8,colframe=ctempo,boxrule=1pt,
    left=6pt,right=6pt,top=2pt,bottom=2pt]
  \textcolor{ctempo}{\textbf{$\blacktriangleright$ #1}}%
  \hfill\textcolor{ctempo}{\small\textit{#2}}
  \end{tcolorbox}\vspace{3pt}}

\titleformat{\section}{\normalsize\bfseries\color{cfala}}{}{0pt}{}
  [{\color{cfala}\titlerule[0.8pt]}]
\titlespacing{\section}{0pt}{8pt}{4pt}

\pagestyle{fancy}\fancyhf{}
\lhead{\textbf{TR2 — Entrega 1} \textbar\ Roteiro 5 min}
\rhead{\textbf{Adaptive HTTP Streaming}}
\cfoot{\thepage}
\renewcommand{\headrulewidth}{0.4pt}

% ─────────────────────────────────────────────────────────────────────────────
\begin{document}

\begin{tcolorbox}[colback=cfala!10,colframe=cfala,boxrule=2pt,arc=4pt,
  left=10pt,right=10pt,top=8pt,bottom=8pt]
\begin{center}
  {\LARGE\bfseries\color{cfala} Roteiro — Entrega 1}\\[4pt]
  {\large Rate-Based ABR $\cdot$ Cliente funcional $\cdot$ Demo ao vivo}\\[2pt]
  \textbf{5 minutos} \quad|\quad
  \small\textcolor{cfala}{FALA = dizer}\quad
  \textcolor{cacao}{AÇÃO = fazer na tela}\quad
  \textcolor{catencao}{ATENÇÃO = cuidado}
\end{center}
\end{tcolorbox}

\vspace{4pt}

% ══════════════════════════════════════════════════════════════════════════════
\section{1. Abertura}
% ══════════════════════════════════════════════════════════════════════════════

\tempo{0:00 – 0:40}{Contexto rápido}

\begin{acaobox}
Antes de começar já ter aberto:\\
\cmd{Terminal} na pasta \cmd{client/} \quad|\quad
\cmd{abr.py} no editor (linha 18) \quad|\quad
\cmd{Wireshark} capturando em \cmd{wlp1s0}
\end{acaobox}

\vspace{3pt}

\begin{falabox}
``Vamos mostrar o cliente de streaming adaptativo da Entrega 1 rodando contra
o servidor real. O cliente baixa o manifest, mede a vazão de cada segmento,
escolhe a qualidade pelo algoritmo Rate-Based ABR e grava tudo em CSV com
gráficos automáticos.''
\end{falabox}

% ══════════════════════════════════════════════════════════════════════════════
\section{2. Código — ABR e Buffer (rápido)}
% ══════════════════════════════════════════════════════════════════════════════

\tempo{0:40 – 1:30}{Mostrar o código}

\begin{acaobox}
Ir para o editor $\to$ \cmd{abr.py} linha 52 — destacar a linha \cmd{safety\_limit = throughput\_kbps * self.SAFETY\_FACTOR}\\[4pt]
Depois ir para \cmd{buffer\_manager.py} linha 45 (\cmd{add\_segment}) e linha 65 (\cmd{consume})
\end{acaobox}

\vspace{3pt}

\begin{falabox}
``O ABR calcula \texttt{safety\_limit = vazão × 0,85} e escolhe a maior qualidade
cujo bitrate caiba nesse limite — os 15\% de margem evitam rebuffering em
flutuações de rede.

O \texttt{BufferManager} faz \texttt{buffer += 2s} a cada segmento recebido e
\texttt{buffer -= tempo\_de\_download} pelo tempo que o player ficou esperando.
Se zerar, rebuffering. Se $\geq$ 2s, o campo \texttt{buffer\_can\_play} é 1.''
\end{falabox}

% ══════════════════════════════════════════════════════════════════════════════
\section{3. Demo ao vivo — cliente + Wireshark}
% ══════════════════════════════════════════════════════════════════════════════

\tempo{1:30 – 4:00}{Demo principal}

\begin{alertabox}
Wireshark já deve estar rodando. Aplicar o filtro agora se ainda não aplicou:\\
\cmd{ip.addr == 137.131.178.229}
\end{alertabox}

\vspace{3pt}

\begin{acaobox}
Rodar o cliente no terminal:\\[4pt]
\cmd{python3 main.py}\\[4pt]
Aguardar o seg 001 aparecer e começar a falar.
\end{acaobox}

\vspace{3pt}

\begin{falabox}
``O manifest foi baixado — no Wireshark viram o SYN e o GET /manifest.

\acaoinline{Apontar para a linha seg 001 no terminal}

Segmento 001 saiu em 240p, cold start sem histórico ainda. A partir do 002 o
ABR já tem uma medição e sobe a qualidade. Olhem as colunas: \texttt{vazao}
é a banda medida, \texttt{buf} é o buffer crescendo em segundos — ele sobe
porque baixamos mais rápido do que a taxa de reprodução — e \texttt{dl} é o
tempo real de cada download.

\acaoinline{Alternar para o Wireshark}

No Wireshark cada rajada de pacotes TCP é um segmento. O intervalo entre
rajadas é o processamento do cliente. Quando o jitter no terminal sobe, dá
para ver o espaçamento irregular entre os pacotes aqui.

\acaoinline{Voltar para o terminal e deixar terminar}

A qualidade estabiliza em 480p porque com o servidor em 2 Mbps e o safety
factor, a vazão média de 1.100 kbps fica abaixo do limiar de 1.500 kbps da
720p. Deficiência do Rate-Based puro que vamos resolver na Entrega 2.''
\end{falabox}

% ══════════════════════════════════════════════════════════════════════════════
\section{4. CSV e gráfico}
% ══════════════════════════════════════════════════════════════════════════════

\tempo{4:00 – 4:45}{Evidências}

\begin{acaobox}
\cmd{cat ../logs/\$(ls -t ../logs/ | head -1) | column -t -s, | head -8}\\[4pt]
\cmd{xdg-open ../graphs/throughput\_timeline.png \&}
\end{acaobox}

\vspace{3pt}

\begin{falabox}
``CSV com as 14 colunas da especificação — \texttt{buffer\_can\_play} em 1 em
todos os segmentos, sem nenhum rebuffering. O gráfico mostra a vazão e a
qualidade selecionada: correlação direta, e o buffer subindo continuamente.
Esse é o baseline da Entrega 1.''
\end{falabox}

% ══════════════════════════════════════════════════════════════════════════════
\section{5. Encerramento}
% ══════════════════════════════════════════════════════════════════════════════

\tempo{4:45 – 5:00}{Fechar}

\begin{falabox}
``Pronto. Cliente funcional, CSV gerado, gráficos automáticos. Alguma pergunta?''
\end{falabox}

\vspace{6pt}
\noindent\rule{\textwidth}{0.8pt}

% ══════════════════════════════════════════════════════════════════════════════
\section*{Respostas rápidas}
% ══════════════════════════════════════════════════════════════════════════════

\begin{tcolorbox}[colback=cinzabg,colframe=gray!50,boxrule=0.8pt,arc=3pt,
  left=8pt,right=8pt,top=5pt,bottom=5pt]
\begin{tabularx}{\textwidth}{>{\bfseries}p{4.5cm}X}
\toprule
Pergunta & Resposta \\
\midrule
Onde o ABR decide? &
  \cmd{abr.py} linha 52 — \cmd{safety\_limit = throughput * 0.85}, itera qualidades ordenadas por bitrate. \\
\addlinespace
Como o buffer é calculado? &
  \cmd{buffer += 2s} (segmento recebido) $-$ \cmd{download\_time\_s} (tempo de espera). \\
\addlinespace
O que é \cmd{buffer\_can\_play}? &
  1 se \cmd{buffer $\geq$ 2s}. Indica reprodução contínua sem interrupção. \\
\addlinespace
Por que parou em 480p? &
  Banda 2 Mbps $\times$ 0,85 $\approx$ 935 kbps efetivos $<$ 1.500 kbps da 720p. \\
\addlinespace
O que é jitter EWMA? &
  Média móvel exponencial (alpha=0,2) da variação entre segmentos consecutivos. \\
\bottomrule
\end{tabularx}
\end{tcolorbox}

\vspace{4pt}
\begin{tcolorbox}[colback=bgatencao,colframe=catencao,boxrule=0.8pt,arc=3pt,
  left=8pt,right=8pt,top=4pt,bottom=4pt]
\small
\textbf{Checklist pré-apresentação:}\quad
\cmd{curl http://137.131.178.229:8080/health}\quad$\cdot$\quad
Wireshark em \cmd{wlp1s0}\quad$\cdot$\quad
Teste rápido \cmd{python3 main.py}\quad$\cdot$\quad
Gráficos de backup em \cmd{graphs/}
\end{tcolorbox}

\end{document}
